% ---------------------------------------------------------------------------
% Eurographics 2027 full-paper submission (Computer Graphics Forum special issue)
% Preamble mirrors eg-style/egPublStyle-cgf/EGauthorGuidelines-cgf-sub.tex.
% TODO: after logging in to SRMv2, diff against the EG_2027 package's
%       EGauthorGuidelines-conf-sub.tex (see SUBMISSION_NOTES.md) and adopt
%       its mode switch if it differs from \JournalSubmission.
% ---------------------------------------------------------------------------
\documentclass{egpubl}

% --- for  CGF Journal
\JournalSubmission    % submission (anonymous) mode
% \JournalPaper       % final version

% !! *please* don't change anything above
% !! unless you REALLY know what you are doing
% ------------------------------------------------------------------------
\usepackage[T1]{fontenc}
\usepackage{dfadobe}

\usepackage{cite}  % comment out for biblatex with backend=biber
% ---------------------------
%\biberVersion
\BibtexOrBiblatex
% ---------------------------
\electronicVersion
\PrintedOrElectronic
\ifpdf \usepackage[pdftex]{graphicx} \pdfcompresslevel=9
\else \usepackage[dvips]{graphicx} \fi

\usepackage{egweblnk}
% end of EG prologue
% ------------------------------------------------------------------------

\usepackage{booktabs}
\usepackage{multirow}
\usepackage{amsmath,amssymb}
\usepackage{xcolor}

% Working-draft helpers. Remove before submission.
\newcommand{\todo}[1]{\textcolor{red}{[TODO: #1]}}

% ------------------------------------------------------------------------
\title[Synthetic LiDAR-Camera Data for Human Pose and Shape Estimation]%
      {Synthetic LiDAR-Camera Data for Human Pose and Shape Estimation}

% Anonymous submission: use the SRMv2 submission ID instead of author names.
\author[SUBMISSION ID]{SUBMISSION ID}

% ------------------------------------------------------------------------
\begin{document}

% TODO: teaser figure -- BEDLAM RGB frame, simulated LiDAR sweep, fused SMPL
% estimate with 3D box, and a real SLOPER4D/Waymo result side by side.
% \teaser{
%   \includegraphics[width=0.9\linewidth]{figures/teaser}
%   \centering
%   \caption{TODO}
%   \label{fig:teaser}
% }

\maketitle

% ------------------------------------------------------------------------
\begin{abstract}
% TODO (1-2 sentences per point): the problem (real LiDAR-camera SMPL data
% is scarce and rarely has metric 3D placement); the dataset (LiDAR
% simulated from BEDLAM depth with Ouster-style beam presets and randomised
% sensor placement); the model (selective-attention fusion that routes 3D
% cues to LiDAR and semantic cues to the camera); the result (synthetic
% pre-training improves SMPL and 3D placement on SLOPER4D, LiDARHuman26M
% and Waymo).
\todo{Write abstract.}
\end{abstract}

% ------------------------------------------------------------------------
\section{Introduction}
\label{sec:intro}
% TODO: motivate LiDAR-camera human mesh recovery (metric depth, outdoor,
% robotics/AV) and the lack of large real datasets with SMPL labels.
% TODO: why synthetic: BEDLAM already provides depth, SMPL-X and camera
% intrinsics, so a physically plausible LiDAR sweep can be derived per frame.
% Contributions (the red thread: data -> model -> evidence, one addition
% per ablation rung; see docs/story.md).
Our contributions are threefold:
\begin{itemize}
  \item \textbf{A LiDAR-enhanced synthetic dataset for SMPL pose and
    shape.} We derive physically plausible LiDAR sweeps from the depth of
    a large-scale synthetic image dataset, making this the first synthetic 
    dataset extension to hold multi-modal SMPL samples for training of fused 
    camera-LiDAR SMPL estimation.
  \item \textbf{A selective-attention fusion model.} Per-query gates route
    semantic joint groups (head, arms, hands, legs, feet) to the image and
    metric quantities (body centre, placement, orientation, shape) to the
    LiDAR, on a frozen image backbone.
  \item \textbf{Extensive experiments and evaluation on two challenging
    in-the-wild datasets.} On Waymo and SLOPER4D we show that our approach
    improves pose, shape and camera-frame 3D placement over image-only and
    LiDAR-only estimation, and ablate one ingredient at a time: fusion,
    data mixture, and synthesis choices. As on any calibrated sensor rig,
    the camera intrinsics are given, so unlike most SMPL estimators we do
    not have to infer them.
\end{itemize}
\todo{Write introduction.}

% ------------------------------------------------------------------------
\section{Related Work}
\label{sec:related}

\subsection{Image-based human mesh recovery}
% TODO: HMR lineage to ViT-based models; TokenHMR~\cite{dwivedi2024tokenhmr}
% and CameraHMR~\cite{patel2025camerahmr}; SAM 3D Body~\cite{sam3dbody2025}.
% Note: none recover metric camera-frame translation reliably without depth.
\todo{Image HMR paragraph.}

\subsection{LiDAR-based human mesh recovery}
% TODO: LiDARCap / LiDARHuman26M~\cite{li2022lidarcap},
% LiDAR-HMR~\cite{fan2023lidarhmr}, SLOPER4D baselines~\cite{dai2023sloper4d};
% sparse points give metric position but weak pose/shape detail.
\todo{LiDAR HMR paragraph.}

\subsection{LiDAR-camera fusion}
% TODO: fusion for 3D detection (Waymo-style pipelines) and for humans
% (LIF-Net~\cite{lifnet}); contrast early/late fusion with our per-query
% modality gating.
\todo{Fusion paragraph.}

\subsection{Synthetic data for human pose}
% TODO: BEDLAM~\cite{black2023bedlam} and its use for image HMR; no prior
% work derives LiDAR from it; domain gap discussion.
\todo{Synthetic data paragraph.}

% ------------------------------------------------------------------------
\section{Synthetic Dataset Generation}
\label{sec:dataset}

\subsection{BEDLAM source and ground truth}
% TODO: which BEDLAM sequences/cameras are used, the SMPL(-X) bodies,
% depth maps and intrinsics; how we convert SMPL-X labels to SMPL and
% express translation in the camera frame; per-person 3D boxes from the
% posed mesh.
\todo{Describe BEDLAM inputs and labels.}

\subsection{LiDAR simulation}
% TODO: Ouster-style presets as (channels x steps per revolution), the
% vertical FoV per family, and the azimuth window that overlaps the camera
% frustum. Sensor placement: origin sampled in a ball around the camera
% centre with an orientation that keeps the subject in view; ray casting
% against the BEDLAM depth map with range/intensity/return fields.
\todo{Describe beam model, sampling, and depth ray casting.}

\subsection{Realism}
% TODO: rolling shutter coupled to ego speed (per-column timestamps and
% motion of the sensor during one revolution); occlusion augmentation
% (dropouts, self/scene occluders); noise, range-dependent dropout and
% intensity augmentation.
\todo{Describe rolling shutter, ego speed, occlusions, augmentations.}

\subsection{Statistics}
% TODO: number of frames / subjects / sweeps, distribution of range,
% points-per-person, sensor offsets; comparison table against SLOPER4D,
% LiDARHuman26M and Waymo.
\todo{Dataset statistics and comparison table.}

\begin{table*}[t]
  \centering
  \caption{Dataset comparison. \todo{fill}}
  \label{tab:dataset_stats}
  \begin{tabular}{@{}lrrrrl@{}}
    \toprule
    Dataset & Frames & Subjects & Points/person & Range (m) & SMPL GT \\
    \midrule
    SLOPER4D~\cite{dai2023sloper4d}  & -- & -- & -- & -- & -- \\
    LiDARHuman26M~\cite{li2022lidarcap} & -- & -- & -- & -- & -- \\
    Waymo~\cite{sun2020waymo}        & -- & -- & -- & -- & -- \\
    Ours (synthetic)                  & -- & -- & -- & -- & -- \\
    \bottomrule
  \end{tabular}
\end{table*}

% ------------------------------------------------------------------------
\section{Selective-Attention Fusion}
\label{sec:method}

\subsection{Backbones}
% TODO: ViT-H image encoder (patch tokens + camera intrinsics conditioning);
% point tokenizer (voxel/kNN grouping into point tokens with positional
% encodings in the camera frame).
\todo{Describe image and point backbones.}

\subsection{Joint-group queries and modality gates}
% TODO: a small set of learned queries per joint group (torso, limbs, head,
% global); each query attends to image and LiDAR tokens through separate
% cross-attention branches; a per-query gate mixes the two so that 3D
% placement is driven by LiDAR and semantics by the image.
\todo{Describe queries, cross-attention routing and gating.}

\subsection{Outputs}
% TODO: SMPL pose/shape heads, camera-frame translation head, and 3D box
% (centre, size, yaw) head; how the box is tied to the SMPL mesh.
\todo{Describe output heads.}

\subsection{Losses}
% TODO: SMPL parameter loss, 3D/2D joint losses, vertex loss, translation
% loss, box regression/IoU loss, and the weighting schedule.
\todo{Describe losses.}

% ------------------------------------------------------------------------
\section{Experiments}
\label{sec:experiments}

\subsection{Datasets and protocol}
% TODO: SLOPER4D, LiDARHuman26M and Waymo splits; how real sequences are
% converted to our format (camera frame, SMPL); training regimes
% synthetic-only / real-only / synthetic+real.
% Metrics: MPJPE, PA-MPJPE, PVE (mm); translation error (cm);
% 3D box AP / mAP (IoU thresholds TODO).
\todo{Describe datasets, splits, metrics.}

\subsection{Baselines}
% TODO: CameraHMR~\cite{patel2025camerahmr}, TokenHMR~\cite{dwivedi2024tokenhmr},
% LiDAR-HMR~\cite{fan2023lidarhmr}, SAM 3D Body~\cite{sam3dbody2025},
% LIF-Net~\cite{lifnet}; how each is evaluated for translation / boxes.
\todo{Describe baselines.}

\subsection{Main results}
% TODO: Table~\ref{tab:main} discussion: synthetic-only transfer, real-only,
% synthetic+real; gains concentrated in translation and box metrics.
\todo{Discuss main results.}

\begin{table*}[t]
  \centering
  \caption{Main results on SLOPER4D, LiDARHuman26M and Waymo. Lower is better for MPJPE, PA-MPJPE, PVE (mm) and translation error (cm); higher is better for 3D box AP / mAP. \todo{fill}}
  \label{tab:main}
  \begin{tabular}{@{}llcrrrrrr@{}}
    \toprule
    Method & Input & Training data & MPJPE & PA-MPJPE & PVE & Trans.\ err. & Box AP & mAP \\
    \midrule
    CameraHMR~\cite{patel2025camerahmr} & RGB   & real      & -- & -- & -- & -- & -- & -- \\
    TokenHMR~\cite{dwivedi2024tokenhmr} & RGB   & real      & -- & -- & -- & -- & -- & -- \\
    SAM 3D Body~\cite{sam3dbody2025}    & RGB   & real      & -- & -- & -- & -- & -- & -- \\
    LiDAR-HMR~\cite{fan2023lidarhmr}    & LiDAR & real      & -- & -- & -- & -- & -- & -- \\
    LIF-Net~\cite{lifnet}               & RGB+LiDAR & real  & -- & -- & -- & -- & -- & -- \\
    \midrule
    Ours & RGB+LiDAR & synthetic only      & -- & -- & -- & -- & -- & -- \\
    Ours & RGB+LiDAR & real only           & -- & -- & -- & -- & -- & -- \\
    Ours & RGB+LiDAR & synthetic + real    & -- & -- & -- & -- & -- & -- \\
    \bottomrule
  \end{tabular}
\end{table*}

% ------------------------------------------------------------------------
\section{Ablation Studies}
\label{sec:ablation}

\subsection{Selective attention: pose accuracy}
% TODO: replace gated routing with concatenation / early fusion / late
% fusion; report MPJPE, PA-MPJPE, PVE.
\todo{Discuss Table~\ref{tab:abl_attention_pose}.}

\begin{table}[t]
  \centering
  \caption{Fusion variants: pose accuracy. \todo{fill}}
  \label{tab:abl_attention_pose}
  \begin{tabular}{@{}lrrr@{}}
    \toprule
    Fusion & MPJPE & PA-MPJPE & PVE \\
    \midrule
    Image only            & -- & -- & -- \\
    LiDAR only            & -- & -- & -- \\
    Early fusion (concat) & -- & -- & -- \\
    Late fusion           & -- & -- & -- \\
    Selective attention (ours) & -- & -- & -- \\
    \bottomrule
  \end{tabular}
\end{table}

\subsection{Selective attention: LiDAR impact on 3D box and mAP}
% TODO: same variants, translation error and box AP/mAP; show that gating
% lets LiDAR own placement without hurting pose.
\todo{Discuss Table~\ref{tab:abl_attention_box}.}

\begin{table}[t]
  \centering
  \caption{Fusion variants: 3D placement. \todo{fill}}
  \label{tab:abl_attention_box}
  \begin{tabular}{@{}lrrr@{}}
    \toprule
    Fusion & Trans.\ err. & Box AP & mAP \\
    \midrule
    Image only            & -- & -- & -- \\
    LiDAR only            & -- & -- & -- \\
    Early fusion (concat) & -- & -- & -- \\
    Late fusion           & -- & -- & -- \\
    Selective attention (ours) & -- & -- & -- \\
    \bottomrule
  \end{tabular}
\end{table}

\subsection{Sensor extrinsics: fixed rig vs.\ random placement}
% Decided plan (docs/ablations.md): camera as is, LiDAR fixed to the target
% rig (Waymo / SLOPER4D calibration tree) vs. random placement in a ball of
% radius 0.25 m vs. 1.0 m around the camera; evaluate on Waymo and SLOPER4D.
% Note: both target rigs have 6-10 cm camera-LiDAR baselines, so the sweep
% tells how much placement variance synthesis needs.
We ask whether synthesising for the target vehicle's calibration tree pays off:
\begin{itemize}
  \item Rig-fixed LiDAR pose vs.\ random ball around the camera.
  \item Does the calibration tree improve accuracy on that rig?
  \item Does it speed up convergence, at what transfer cost?
\end{itemize}
\todo{Discuss Table~\ref{tab:abl_extrinsics}.}

\begin{table*}[t]
  \centering
  \caption{Sensor placement during synthesis. \todo{fill}}
  \label{tab:abl_extrinsics}
  \begin{tabular}{@{}llrrrrrr@{}}
    \toprule
    Placement & Test set & MPJPE & PA-MPJPE & PVE & Trans.\ err. & Box AP & mAP \\
    \midrule
    Fixed to rig & Waymo    & -- & -- & -- & -- & -- & -- \\
    Random ball  & Waymo    & -- & -- & -- & -- & -- & -- \\
    Fixed to rig & SLOPER4D & -- & -- & -- & -- & -- & -- \\
    Random ball  & SLOPER4D & -- & -- & -- & -- & -- & -- \\
    \bottomrule
  \end{tabular}
\end{table*}

\subsection{LiDAR resolution: channels and steps per revolution}
% TODO: sweep Ouster presets (e.g. 32/64/128 channels x 512/1024/2048 steps);
% accuracy vs. points per person.
\todo{Discuss Table~\ref{tab:abl_resolution}.}

\begin{table*}[t]
  \centering
  \caption{LiDAR resolution. \todo{fill}}
  \label{tab:abl_resolution}
  \begin{tabular}{@{}rrrrrrrr@{}}
    \toprule
    Channels & Steps/rev & MPJPE & PA-MPJPE & PVE & Trans.\ err. & Box AP & mAP \\
    \midrule
    32  & 512  & -- & -- & -- & -- & -- & -- \\
    32  & 1024 & -- & -- & -- & -- & -- & -- \\
    64  & 1024 & -- & -- & -- & -- & -- & -- \\
    128 & 1024 & -- & -- & -- & -- & -- & -- \\
    128 & 2048 & -- & -- & -- & -- & -- & -- \\
    \bottomrule
  \end{tabular}
\end{table*}

\subsection{Realism: augmentations and occlusions}
% Decided plan: augmentations (image + point) and occlusions are ablated
% after hand-in. Rolling shutter with ego speed is a synthesis option
% (Section 3.3) but not an ablation: Waymo's top-LiDAR points are
% motion-compensated by the official tooling, so the effect cannot show on
% the primary target.
\todo{Discuss Table~\ref{tab:abl_realism}.}

\begin{table*}[t]
  \centering
  \caption{Realism components in synthesis. \todo{fill}}
  \label{tab:abl_realism}
  \begin{tabular}{@{}lrrrrrr@{}}
    \toprule
    Configuration & MPJPE & PA-MPJPE & PVE & Trans.\ err. & Box AP & mAP \\
    \midrule
    Full (ours)                        & -- & -- & -- & -- & -- & -- \\
    without augmentations              & -- & -- & -- & -- & -- & -- \\
    without occlusions                 & -- & -- & -- & -- & -- & -- \\
    \bottomrule
  \end{tabular}
\end{table*}

% ------------------------------------------------------------------------
\section{Limitations}
\label{sec:limitations}
% TODO: depth-derived LiDAR misses multi-return and material effects;
% BEDLAM scenes are indoor/urban-limited; SMPL (not SMPL-X) outputs;
% single-person focus.
\todo{Write limitations.}

% ------------------------------------------------------------------------
\section{Conclusion}
\label{sec:conclusion}
% TODO: restate the three contributions and the headline gain in metric
% placement; future work on multi-person and temporal models.
\todo{Write conclusion.}

% ------------------------------------------------------------------------
\bibliographystyle{eg-alpha-doi}
\bibliography{refs}

\end{document}
